\section{Introducción: cómo las capacidades de software de la IA se extienden al hardware}

Esta sección establece primero los objetivos de aprendizaje de todo el episodio. Este episodio sobre Rokid no es una simple historia oral de emprendimiento, sino un caso de estudio sobre «el punto de acceso del hardware de IA de próxima generación». Zhu Mingming (祝铭明) pasó por cuatro ciclos, el internet móvil, las bocinas con IA, los lentes de AR y los modelos grandes: su primera empresa fue adquirida por Alibaba, participó en M Lab dentro de Alibaba y emprendió por segunda vez con Rokid. Al leer este episodio, lo importante no es recordar en qué año hubo una ronda de financiamiento, sino entender cómo un emprendedor de hardware evalúa el punto de acceso, la plataforma, la cadena de suministro, la organización y la competencia con los gigantes.

Zhang Xiaojun (张小珺) dice al inicio que, a medida que las capacidades de software de la IA se extienden al hardware, además de la inteligencia corporizada, los lentes inteligentes podrían ser otra industria beneficiada. Este juicio sitúa el EP104 dentro de la serie de IA/internet de Zhang Xiaojun: no es una entrevista de electrónica de consumo pura, sino una discusión sobre cómo las capacidades de la IA pasan de la nube, el teléfono celular y los productos de voz a un dispositivo personal always-on.

\begin{importantbox}{Tesis central de este episodio}
La oportunidad de los lentes inteligentes no consiste solo en «meter la IA en unos lentes», sino en redefinir el punto de acceso personal a la información. El juicio de Rokid es este: la IA necesita un soporte disponible en cualquier momento y lugar, que entregue la información de forma directa y que tenga capacidad de visualización; el teléfono celular seguirá existiendo, pero las interacciones fragmentadas podrían trasladarse poco a poco a los lentes.
\end{importantbox}

\begin{figure}[H]
\centering
\includegraphics[width=0.96\textwidth]{ai-to-hardware-spillover.png}
\caption{Las capacidades de software de la IA se extienden al hardware: del modelo en la nube a la inteligencia portátil y la inteligencia corporizada. Diagrama conceptual propio, elaborado a partir de la conversación en 00:01:12--00:02:00 y 01:11:52--01:13:05.}
\end{figure}

\begin{knowledgebox}{Lectura de la figura: la inteligencia corporizada y la inteligencia portátil son dos rutas hacia el hardware}
En la figura, AI Model se extiende hacia Device y se divide en dos categorías: Wearable y Embodied. La inteligencia corporizada pone el énfasis en que los robots actúen en el mundo físico; la inteligencia portátil, en que las personas puedan recurrir a la IA en cualquier momento. Los lentes inteligentes pertenecen a la segunda: lo decisivo en ellos es la eficiencia como punto de acceso, la visualización y la comodidad al llevarlos puestos, no la ejecución mecánica.
\end{knowledgebox}

\begin{knowledgebox}{Nota sobre la estrategia visual}
Este video es una entrevista con encuadre fijo, sin slides, pizarrón ni demostraciones de producto. El texto solo usa la portada para identificar la fuente; todas las imágenes del cuerpo son diagramas conceptuales propios que explican la cronología del emprendimiento, el giro hacia la IA y la AR, el «bosque oscuro» del hardware, las etapas de los lentes inteligentes y las diferencias en la definición de producto entre China y Estados Unidos.
\end{knowledgebox}

\subsection{Ruta de lectura: la historia del emprendimiento como lección sobre cómo evaluar puntos de acceso}

Esta sección propone además una forma de leer. La entrevista contiene una gran cantidad de nombres de personas, rondas de financiamiento, productos y años; si se memorizan uno por uno, es fácil leerla como una «historia del fundador». Una lectura más valiosa es devolver cada episodio de la historia a un marco de juicio: cuando un nuevo hardware busca convertirse en punto de acceso, qué condiciones debe cumplir exactamente; y cuando un producto solo está aprovechando una tendencia, qué señales indican que es un producto y no una plataforma.

\begin{knowledgebox}{Nota de clase: este episodio no es una evaluación de especificaciones de lentes}
Aquello a lo que el ponente vuelve una y otra vez no es una especificación concreta, sino el juicio abstracto sobre el «punto de acceso»: tiempo de uso, amplitud de entrada y salida, frecuencia de las interacciones fragmentadas, naturalidad al llevarlos puestos, capacidad de visualización, posición en el ecosistema y capacidad de entrega de la cadena de suministro. Las especificaciones solo están al servicio de estos juicios; no pueden, a la inversa, sustituirlos.
\end{knowledgebox}

\noindent\begin{tabularx}{\textwidth}{p{0.18\textwidth}p{0.36\textwidth}X}
\toprule
Pregunta de lectura & Material en la entrevista & Juicio que se debe formar \\
\midrule
¿De dónde surge el punto de acceso? & iPhone, sistemas operativos móviles, bocinas con IA, lentes inteligentes & El cambio de punto de acceso suele ocurrir cuando la forma del hardware y las capacidades del software cambian al mismo tiempo. \\
¿Cuándo se consolida una plataforma? & Dos horas de uso diario promedio, desarrolladores, usuarios empresariales, ecosistema de contenido & Una plataforma no existe porque la empresa lo proclame, sino que se forma por la combinación de la frecuencia de uso y la oferta del ecosistema. \\
¿Por qué hace falta visualización? & Traducción, avisos, consulta rápida de información, resultados intermedios de baja confianza & La visualización convierte la salida de la IA, de voz en serie, en una interfaz visible y corregible. \\
¿Cómo sobrevive una empresa emergente? & Los cuatro «no», las ventanas de oportunidad, «acumular grano y levantar murallas altas», hacer muchos aliados & En la etapa temprana, la ventana se aprovecha con velocidad; la competencia a largo plazo depende del producto, el ecosistema y la paciencia del capital. \\
\bottomrule
\end{tabularx}

\subsection{Resumen de la sección}

El hilo conductor del EP104 es este: cuando las capacidades de la IA maduran, los puntos de acceso de hardware se reordenan. El valor de Rokid como caso de estudio está en que, con el costo de más de diez años de emprendimiento en hardware, explica por qué los lentes inteligentes podrían convertirse en el punto de acceso de la IA portátil. Las secciones siguientes explicarán primero de dónde proviene este juicio sobre el punto de acceso y, después, por qué se abandonaron las bocinas y se persistió en los lentes, cómo llegó la competencia de los gigantes y por qué al emprendimiento en hardware se le llama un «bosque oscuro».
